\section{Current compositionality benchmarks and their biases}

A majority of existing compositionality benchmarks for vision-language models formulate the evaluation task as image-to-text retrieval~\cite{zhao2022vl,yuksekgonul2023when,ma2022crepe}. We focus on these benchmarks and discuss others~\cite{thrush2022winoground,ray2023cola} in Appendix~\ref{app:bms}.
Given an image, the model is probed to select text that correctly describes the image from a pool of candidates.
Unlike standard retrieval tasks where the negative (incorrect) candidates differ a lot from the \emph{positive} (correct) text, compositionality benchmarks intentionally design \emph{hard negative} texts that differ minimally from the positive text, in order to test whether the model understands the fine-grained atomic concepts that compose the scene.
% This framework, illustrated in Figure~\ref{fig:hn_examples}, encompasses most recent compositionality benchmarks albeit their different specific forms of hard negatives considered.

% \begin{figure}
%     \centering
%     \includegraphics[width=\linewidth]{imgs/hn_examples.png}
%     \caption{Image-to-text retrieval for evaluating VLMs' compositionality.
%     \aniline{1. The image needs to be saved as a vector graphic so that the text scales up ok. Preferably as a PDF. Ping me if you need help. 2. This isnt just your method, but rather what is typically done right ? That needs to be made clear in the image.} \rk{Image size can be reduced. General rule of thumb: all text in figures and tables should be at least as large as the main text.} \rk{Also, captions should be written so that the first sentence is always the main takeaway from the figure/table. Write captions so that if a reader only reads the captions of your paper, they should understand all the main points.\zx{figure 1 has been converted into table 1 for readability}}}
%     \label{fig:hn_examples}
% \end{figure}

% % \begin{table*}[t]
%      \centering
%      \caption{Following existing compositionality benchmarks for vision-language models, \method adopts image-to-text retrieval for evaluation, where the query is an image and the retrieval set consists of one positive caption plus (at least) one composition-aware hard negative text. }
%      \resizebox{\textwidth}{!}{
%      {\renewcommand{\arraystretch}{1.15}
% \begin{tabular}{lllll}
% \small
% \textbf{Image}     & \multicolumn{4}{l}{\textbf{Text retrieval set}}                                                                                                              \\ 
%          \multirow{8}{*}{\includegraphics[width=0.3\textwidth]{imgs/img_example.jpg}}          & \multicolumn{3}{l}{Positive}                                   & \positive{Set of toy animals sitting in front of a red wooden wagon.}                         \\\cline{2-5}
%                    &                                 &                           & Object    & \negative{Balloons} sitting in front of a red wooden wagon.                                   \\
%                    &                                 &                           & Attribute & Set of \negative{real} animals sitting in front of a red wooden wagon. \\
%                    &                                 & \multirow{-3}{*}{Replace} & Relation  & Set of toy animals \negative{inside} a red wooden wagon.                                      \\\cline{3-5}
%                    &                                 &                           & Object    & \negative{Red wooden wagon} sitting in front of a \negative{set of toy animals}.                         \\
%                    &                                 & \multirow{-2}{*}{Swap}    & Attribute & Set of \negative{wooden} animals sitting in front of a red \negative{toy} wagon.                         \\\cline{3-5}
%                    &                                 &                           & Object    & \negative{Flowers} and a set of toy animals are sitting in front of a red wooden wagon.       \\
% \multirow{-8}{*}{} & \multirow{-7}{*}{Hard Negative} & \multirow{-2}{*}{Add}     & Attribute & A \negative{mechanical} set of toy animals sitting in front of a red wooden wagon. \\\cline{2-5}          
% \end{tabular}}}
% \label{tab:img2txt_retrieval}
% \end{table*}
  



\textbf{Existing hard negative generation process introduces undesirable biases.}
% Existing benchmarks typically generate hard negative texts through rule-based programmatic procedures~\cite{zhao2022vl,yuksekgonul2023when,ma2022crepe} from dense scene annotations in Visual Genome~\cite{krishna2017visual} and MSCOCO~\cite{lin2014microsoft}. 
% % Given an image-text pair, the hard negatives are generated by applying various types of transformations to the original text. 
% These procedures produces hard negatives by replacing a word of specific type (an object, attribute, or relation) in the original text, by swapping two words, or by shuffling the word order (Figure~\ref{fig:hn_examples}). 
Existing benchmarks generate hard negative texts through rule-based programmatic procedures~\cite{zhao2022vl,yuksekgonul2023when,ma2022crepe}, which produce hard negatives by replacing a word of specific type (an object, attribute, or relation) in the original text, by swapping two words, or by shuffling the word order. 
We find that such procedures introduce unintentional biases in the generated hard negatives (see Table~\ref{tab:bad_hn_examples});
specifically, we observe two major types of undesirable artifacts: (1) \emph{nonsensical} artifacts, and (2) \emph{non-fluent} artifacts.
We then utilize Vera~\cite{liu2023vera}, a plausibility estimation model, to characterize the nonsensical bias. 
To capture the non-fluent bias, we leverage a grammar-check model~\cite{morris2020textattack} that assigns high scores to grammatically correct texts.
We find that Vera and the grammar model assign higher scores to positive texts, suggesting that many hard negatives are nonsensical and not fluent (Figure~\ref{fig:vera_gap_existing}).


% \begin{table*}[!t]
%   \centering
%   \small
%   \caption{A table with example HNs that don't make logical sense or have grammatical errors.}
%   \scalebox{0.8}{
% \begin{tabular}{lll}
% Dataset      & Nonsense HNs                                                                                                                                                                                         & Non-fluent HNs / HNs with grammar errors                                                                                                                                                                                        \\ \hline \hline
% % Winoground   & \begin{tabular}[c]{@{}l@{}}Clouds in the head\\ The cow is over the moon\\ Shipping a man\end{tabular}                                                                                            & \begin{tabular}[c]{@{}l@{}}Many a table tennis racket, just one ball\\ One body share with two head\\ A box containing many a car\end{tabular}                                                                                  \\ \hline
% % Cola         & \begin{tabular}[c]{@{}l@{}}Black car of red long train\\ Blue sky to the left of stoned white building\\ Sunny building near brown full large massive tall tree\end{tabular}                         & \begin{tabular}[c]{@{}l@{}}Wood empty bench near flat white wood fence\\ Leather baseball black glove to the left of the white uniform\\ Brown walking muddy large dirty elephant...\end{tabular} \\ \hline
% VL Checklist & \begin{tabular}[c]{@{}l@{}}Sheep is hardwood\\ empty zebras\\ the bush speaking in the garden\end{tabular}                                                                                           & \begin{tabular}[c]{@{}l@{}}an man fishing a food from a wrapper using a paw at a open\\ It heaving at a city\\ an grouping subduing at a room access\end{tabular}                                                               \\ \hline
% ARO          & \begin{tabular}[c]{@{}l@{}}The grass is eating the horse\\ A gray bathtub is looking at a white cat...\\ Green ball with a remarkable chair behind a blue scene\end{tabular} & \begin{tabular}[c]{@{}l@{}}At brown cat a in looking a gray dog sitting is and white bathtub\\ scene with remarkable a ball blue a green behind chair\\ Books the looking at people are\end{tabular}                            \\\hline
% CREPE        & \begin{tabular}[c]{@{}l@{}}olives and grape inside a plate\\ ground in a basket on the flowers\\ a hair wearing a necklace, with her lady on a table\end{tabular}                                    & \begin{tabular}[c]{@{}l@{}}A door with panes not in a room; the door has windows\\ right has word another word. there is a words.\\ A shelf with books in something. There is no background\end{tabular}
% \\\bottomrule
                       
% \end{tabular}
% }
% \label{tab:bad_hn_examples}
% \end{table*}


\begin{table*}[t]
  \centering
  \small
  \caption{Existing compositionality benchmarks rely on procedurally-generated hard negatives which often do not make logical sense or are not fluent due to grammatical errors.}
  \scalebox{0.7}{
\begin{tabular}{lll}
\toprule
    Dataset & Nonsensical Hard Negatives  & Non-fluent Hard Negatives  \\
\midrule

CREPE~\cite{ma2022crepe} & Olives and grape inside a plate. & A door with panes not in a room; the door has windows. \\
 & Ground in a basket on the flowers. & Right has word another word. There is a words. \\
 & A hair wearing a necklace, with her lady on a table. &  A shelf with books in something. There is no background. \\  
 
\midrule

ARO~\cite{yuksekgonul2023when} & The grass is eating the horse. & At brown cat a in looking a gray dog sitting is and white bathtub. \\
 & A gray bathtub is looking at a white cat. &  Scene with remarkable a ball blue a green behind chair. \\
 & Green ball with a remarkable chair behind a blue scene. & Books the looking at people are. \\  

\midrule

VL-CheckList~\cite{zhao2022vl} & Sheep is hardwood. & An man fishing a food from a wrapper using a paw at a open.\\
 & Empty zebras. & It heaving at a city. \\
 & The bush speaking in the garden. & An grouping subduing at a room access.
 
\\\bottomrule
                       
\end{tabular}
}
\label{tab:bad_hn_examples}
\vspace{-2ex}
\end{table*}


\begin{figure}[t]
    \centering
    \includegraphics[width=\linewidth]{imgs/score_gap_existing.png}
    \caption{Top row: We define \emph{Vera score gap} as the score difference between the positive and hard negative texts: $\mathrm{Vera}(T^\mathrm{p}) - \mathrm{Vera}(T^\mathrm{n})$. The entire Vera score gap distribution lies on the positive spectrum, indicating that the template-generated hard negative texts usually have low plausibility. Bottom row: Similarly, \emph{Grammar score gap} is defined by: $\mathrm{Grammar}(T^\mathrm{p}) - \mathrm{Grammar}(T^\mathrm{n})$. On grammar score, we also find that the distribution largely rests on the positive side, suggesting that most hard negative texts in existing benchmarks exhibit grammatical errors.
    }
    \label{fig:vera_gap_existing}
\end{figure}

% \begin{figure}[t]
%     \centering
%     \includegraphics[width=\linewidth]{imgs/grammar_gap_existing.png}
%     \caption{We define the \emph{Grammar score gap} for an example by: $\mathrm{Grammar}(x_\mathrm{pos}) - \mathrm{Grammar}(x_\mathrm{hn})$. We find that the distribution largely rests on the positive side, suggesting that most hard negative texts in existing benchmarks exhibit grammatical.
%     \aniline{It will be good if Fig 2 and 3 are right below each other. Or you can merge into 1 figure}}
%     \label{fig:grammar_gap_existing}
% \end{figure}



\textbf{Dataset biases render current compositionality benchmarks ineffective.}
Given the heavily-skewed score gaps, we show that blind models (\ie, Vera and the grammar model) that simply select the higher-scoring texts as positives and admittedly do not possess any vision-language compositionality, can achieve state-of-the-art performances on existing benchmarks.
We compare the the blind models against $17$ pretrained CLIP models from three sources: OpenAI in-house data~\cite{RadfordKHRGASAM21}, LAION~\cite{laion5b}, and Datacomp~\cite{gadre2023datacomp}.
We plot the performances of the blind models and the best-performing CLIP models from each category (Figure~\ref{fig:eval_existing}).
Blind models achieves state-of-the-art performances on $9$ out of $10$ existing benchmark tasks. We provide full evaluation results in Appendix~\ref{app:eval}.
% These results suggest that current benchmarks are hackable and run the risks of overestimating a model's true vision-language compositionality.




\begin{figure}[t]
    \centering
    \includegraphics[width=\linewidth]{imgs/eval_existing.png}
    \caption{Blind commonsense Vera model and Grammar model outperform state-of-the-art CLIP models on nearly \emph{all} existing benchmarks by exploiting the nonsensical and non-fluent artifacts. This suggests that existing benchmarks are hackable and ineffective in measuring compositionality.
    }
    \label{fig:eval_existing}
\end{figure}











